\documentclass{article}
\usepackage{amsmath,amssymb,amsthm,mathtools}
\usepackage{leanify}
\newtheorem{lemma}{Lemma}
\newtheorem{proposition}[lemma]{Proposition}
\title{The language of proofs: assumptions, cases, contradictions and induction}
\begin{document}
\maketitle

A proof makes hypotheses: it assumes something, it distinguishes cases, it supposes the contrary of what it proves.
Leanify reads each as a hypothesis of the steps that are made under it.

\section{Assumptions}

\begin{lemma}\label{lem:assume}
Let $x\in\mathbb{R}$.
Assume that $x>1$.
Then $x^2>x$.
\end{lemma}

\begin{proof}
We have $x-1>0$ and $x>0$.
Hence $x^2-x=x(x-1)>0$.
\end{proof}

\begin{lemma}\label{lem:further}
For every $x\in\mathbb{R}$ and every $y\in\mathbb{R}$, if $0<x<y$, then $x^2<y^2$.
\end{lemma}

\begin{proof}
Let $x\in\mathbb{R}$ and $y\in\mathbb{R}$.
Suppose that $0<x$ and $x<y$.
Then $y-x>0$.
Assume further that $y+x>0$.
Then $y^2-x^2=(y-x)(y+x)>0$.
\end{proof}

\begin{lemma}\label{lem:parity}
Let $n$ be an integer.
Suppose $n$ is even.
Then $n^2$ is even.
\end{lemma}

\begin{lemma}\label{lem:new-variable}
Let $x\in\mathbb{R}$.
Assume $c>0$.
Then $x^2+c>0$.
\end{lemma}

\section{Cases}

\begin{lemma}\label{lem:abs}
For every $x\in\mathbb{R}$, $x\le\lvert x\rvert$.
\end{lemma}

\begin{proof}
Let $x\in\mathbb{R}$.
We distinguish two cases:
\begin{enumerate}
\item \emph{$x\ge0$.} Then $\lvert x\rvert=x$.
\item \emph{Case 2: $x<0$.} Then $\lvert x\rvert=-x$.
Hence $x<0<\lvert x\rvert$.
\end{enumerate}
In both cases, $x\le\lvert x\rvert$.
\end{proof}

\section{Contradictions}

\begin{lemma}\label{lem:contradiction}
For every $x\in\mathbb{R}$, if $x^2<1$, then $x<1$.
\end{lemma}

\begin{proof}
Let $x\in\mathbb{R}$.
Assume that $x^2<1$.
Suppose, for contradiction, that $x\ge1$.
Then $x^2\ge x\ge1$.
This is a contradiction.
Hence $x<1$.
\end{proof}

\begin{lemma}\label{lem:short}
For every $x\in\mathbb{R}$, $x^2\ne-1$.
\end{lemma}

\begin{proof}
Let $x\in\mathbb{R}$.
Suppose that $x^2=-1$.
Then $0\le x^2<0$, a contradiction.
Hence $x^2\ne-1$.
\end{proof}

\section{What is not read}

A second assumption after steps under the first does not say whether it is added or begins another case, and a case in the running proof does not say where it ends.

\begin{lemma}\label{lem:refused}
For every $x\in\mathbb{R}$, $x^2\ge0$.
\end{lemma}

\begin{proof}
Let $x\in\mathbb{R}$.
Assume that $x\ge0$.
Then $x^2=x\cdot x\ge0$.
Assume that $x<0$.
Assume now that $x<0$.
Without loss of generality, assume that $x\ge0$.
\end{proof}

A contradiction refutes the hypothesis made for it, or the one assumption of the proof.
After several assumptions, none of them made for contradiction, it does not say which one it refutes.

\begin{lemma}\label{lem:which}
For every $x\in\mathbb{R}$, if $x>0$, then $x+1>0$.
\end{lemma}

\begin{proof}
We argue by contradiction.
Let $x\in\mathbb{R}$.
Assume that $x>0$.
Assume further that $x+1\le0$.
Then $x\le-1<0$.
This is a contradiction.
\end{proof}

\section{Induction}

A proof by induction states its hypothesis as a sentence: the statement of the result, about the variable that the result introduces.
The base case and the induction step are the items of a list, like other cases.

\begin{lemma}\label{lem:induction}
Let $n$ be a positive integer.
Then $n^2\ge n$.
\end{lemma}

\begin{proof}
We argue by induction on $n$.
\begin{itemize}
\item \emph{Base case: $n=1$.} We have $n^2=1$.
\item \emph{Induction step.} Assume that the statement holds for $n$.
Then $(n+1)^2=n^2+2n+1\ge n+2n+1$.
By the induction hypothesis, $(n+1)^2\ge n+1$.
\end{itemize}
\end{proof}

A statement that binds its variable itself is not a statement about one value of it, and the statement for $n-1$ or for all smaller values is another statement.

\begin{lemma}\label{lem:own-binder}
For every positive integer $n$, $n^2\ge n$.
\end{lemma}

\begin{proof}
Let $n$ be a positive integer.
We argue by induction on $n$.
Assume that the statement holds for $n$.
Assume that the statement holds for $n-1$.
\end{proof}
\end{document}
